\section{한계}
우리가 제안한 다중 에이전트 벤치마크와 프레임워크는 다양한 과제와 평가 지표를 제공하지만, 적용 가능성과 강건성을 높이기 위해서는 다음 영역에 대한 추가 탐구가 필요하다.

\paragraph{Scenario와 Model Coverage 확장.}
현재 우리의 벤치마크는 research co-authoring, Minecraft building, database error analysis, coding collaboration, Werewolf와 bargaining 같은 일부 competitive scenario 등 특정 domain에 초점을 맞춘다. 현실 세계의 다중 에이전트 상호작용 복잡성을 더 잘 포착하기 위해 향후 연구는 open-world environment, 더 풍부한 social cognition이 필요한 시나리오, task-oriented dialogue 같은 application-side task를 포함한 더 다양한 설정을 통합할 수 있다. 모델 측면에서도 본 연구는 전체 스펙트럼을 포괄하지 않는다. 향후 연구에는 DeepSeek model family 같은 다른 최신 모델의 결과가 포함될 수 있다.

\paragraph{Ablation Study 강화.}
현재 분석은 주로 전체 coordination 및 competition performance에 초점을 맞추므로, 특정 component에 대한 더 세밀한 통찰은 충분히 탐구되지 않았다. 향후 실험은 long-term memory, short-term memory, shared memory 같은 서로 다른 memory mechanism과 다양한 multiagent workflow method에 초점을 맞출 수 있다.

\paragraph{Competition Mechanism 발전.}
우리 벤치마크는 competitive task를 포함하지만, multi-party negotiation, repeated strategic play, stochastic element를 수반하는 현실 세계 다중 에이전트 상호작용의 복잡성을 완전히 포착하지는 못한다. 변화하는 환경에서 에이전트가 cooperative role과 adversarial role 사이를 어떻게 전환하는지 조사하는 것은 여전히 유망한 방향이다.

\paragraph{Open-Ended 및 Ill-Defined Task 처리.}
우리 프레임워크의 대부분 과제는 research proposal 완성 또는 database inconsistency 해결처럼 잘 정의된 objective를 포함한다. 그러나 현실 세계 응용에서는 명확한 success criteria가 없는 open-ended 또는 ambiguous context에서 에이전트가 작동해야 하는 경우가 많다. 향후 확장은 다중 에이전트 시스템이 exploratory, non-goal-oriented scenario에 어떻게 적응하는지 탐구할 수 있다.
